\documentclass[10pt,a4paper]{amsart}
\usepackage[margin=19mm]{geometry}
\usepackage{amsmath,amssymb,mathtools}
\usepackage[scr=rsfs,cal=euler]{mathalfa}
\usepackage{xfrac}
\usepackage[unicode,hidelinks,bookmarks=false]{hyperref}
\newcommand{\ie}{i.e.\ }
\newcommand{\cf}{cf.\ }
\newcommand{\resp}{respectively\ }
\newcommand{\Subsection}{\subsection}
\providecommand{\nobreakdash}{\nobreak\hbox{-}}
\setlength{\emergencystretch}{2em}
\multlinegap0pt
\usepackage{amssymb}
\usepackage{algorithm}\usepackage{algpseudocode}
\usepackage[shortlabels]{enumitem}
\usepackage{tikz}
\newtheorem{theorem}{Theorem}
\title{\textbf{Analytic proofs of Andrews-Bachraoui identities related to two-color partitions with evens in one color}}
\author{Gaurab Bardhan, Nipen Saikia}
\date{Source: arXiv:2607.07079v1 (CC BY 4.0)}
\begin{document}
\maketitle

\noindent\emph{Source and licence.} \textbf{Analytic proofs of Andrews-Bachraoui identities related to two-color partitions with evens in one color}. Version arXiv:2607.07079v1 (CC BY 4.0), distributed by arXiv under the Creative Commons Attribution 4.0 International licence, http://creativecommons.org/licenses/by/4.0/, as recorded in the arXiv licence metadata for this exact version and shown on the abstract page https://arxiv.org/abs/2607.07079v1. Full licence notice: \texttt{licenses/arxiv-2607.07079-CC-BY-4.0.txt}.\par\smallskip

\noindent\textit{Editorial scope.} Counted unit: the article's theorem fg3 (source0.tex lines 1289--1306). The article's complete original proof of it is delivered with it (source0.tex lines 1307--1534, one \texttt{proof} environment of 228 lines). No mathematics is written, repaired or re-derived: the statement and the proof are the article's own lines, in the article's own notation, and no retained line is substituted or re-flowed.  Retained uncounted as the source-local context the proof needs: lines 279--281; lines 303--305; lines 309--311; lines 315--317; lines 348--350; lines 384--395; lines 402--404; lines 1071--1083.  Nothing was omitted from any retained span, so the package carries no omission record.  No retained line contains a citation, so the wrapper carries no bibliography and no bibliography entry.  No retained line contains a figure environment, an image-inclusion command or drawing code, so no image is delivered and the package declares no asset dependency.  Editorial formatting only, in addition: an AMS \texttt{amsart} preamble that restates the article's own declarations and macros used by the retained lines, namely the environments \texttt{theorem} on separate counters, together with the standard packages those lines need.  The retained environments print in this document as Theorem 1 in order of appearance, whereas the article prints the same items with the numbering of its own full document.  The complete native source of the article as distributed in the arXiv e-print for this exact version is bundled unchanged as \texttt{sources/204680/source0.tex}.
\begin{equation}\label{psi}
    \psi(q) := f(q, q^3) = \sum_{m=0}^{\infty} q^{m(m+1)/2} = \dfrac{f_2^2}{f_1},
\end{equation}

\begin{equation}\label{gp}
	\sum_{n=0}^\infty r_2(n)q^n=\varphi^2(q),
	\end{equation}

\begin{equation}\label{r_2}
    r_2(n)=4\sum_{\substack{d|n\\d ~odd}}(-1)^{(d-1)/2}=4(d_{1,4}(n)-d_{3,4}(n)),
\end{equation} where, here and throughout the paper, $d_{k,j}(n)$ denotes the number of positive divisors of $n$ such that $d\equiv j\pmod k$.

\begin{equation}\label{r_2.4n+1}
r_2(4n+1)=4\sum_{d\mid 4n+1}\left(\dfrac{-1}{d}\right).
\end{equation}

\begin{equation}\label{B3}
r_{1,2}(4n+1)=2\sum_{d\mid 4n+1}\left(\dfrac{-2}{d}\right).
\end{equation}

   \begin{align}
&\label{P6}\dfrac{1}{f_1^2}
=
\dfrac{f_8^5}{f_2^5f_{16}^2}
+
2q\dfrac{f_4^2f_{16}^2}{f_2^5f_8},\\
&\label{P7}\dfrac{1}{f_1^4}
=
\dfrac{f_4^{14}}{f_2^{14}f_8^4}
+
4q\dfrac{f_4^2f_8^4}{f_2^{10}}.
\end{align}

\begin{equation}\label{2t}
    f_1^{2^t}\equiv f_{2}^{2^{t-1}}\pmod{2^t}. 
\end{equation}   

\begin{align}
\sum_{n=0}^{\infty}F_0(n)q^n
&=
\dfrac{1}{2}
\left(
\dfrac{f_8^5}{f_2^4f_{16}^2}
+
2q\dfrac{f_4^2f_{16}^2}{f_2^4f_8}
+
\dfrac{f_4}{f_2^2}
\right).
\label{A14f0}
\end{align}

\begin{theorem}
For any integer $n\ge 0$ and  $i\in\{0,1\}$, we have 
\begin{equation}\label{fg3}
F_i(4n+1)\equiv
\varepsilon(n)+2\bigl(\mathcal{D}(n)-2\varepsilon(n)\bigr)
\pmod 8.
\end{equation}
Equivalently,
\begin{equation}\label{fg4}
F_i(4n+1)\equiv
\begin{cases}
0 \pmod 8, & \text{if } \varepsilon(n)=0 \text{ and } \mathcal{D}(n)\equiv0\pmod4,\\
4 \pmod 8, & \text{if } \varepsilon(n)=0 \text{ and } \mathcal{D}(n)\equiv2\pmod4,\\
1 \pmod 8, & \text{if } \varepsilon(n)=1 \text{ and } \mathcal{D}(n)\equiv2\pmod4,\\
5 \pmod 8, & \text{if } \varepsilon(n)=1 \text{ and } \mathcal{D}(n)\equiv0\pmod4.
\end{cases}
\end{equation}
\end{theorem}

\begin{proof}
    Extracting the terms involving odd powers of $q$ from \eqref{A14f0}, dividing by $2q,$ and replacing $q^2$ by $q$,  we obtain
\begin{equation}\label{e24}
\sum_{n=0}^{\infty}F_0(2n+1)q^n
=
\dfrac{f_2^2f_8^2}{f_1^4f_4}.
\end{equation}
Employing \eqref{P7} in \eqref{e24}, we obtain
\begin{equation}\label{e25}
\sum_{n=0}^{\infty}F_0(2n+1)q^n
=
\dfrac{f_2^2f_8^2}{f_4}
\left(
\dfrac{f_4^{14}}{f_2^{14}f_8^4}
+
4q\dfrac{f_4^2f_8^4}{f_2^{10}}
\right) 
=
\dfrac{f_4^{13}}{f_2^{12}f_8^2}
+
4q\dfrac{f_4f_8^6}{f_2^8}.
\end{equation}
Extracting the terms involving even powers of $q$ from \eqref{e25} and replacing $q^2$ by $q,$ we obtain
\begin{equation}\label{e26}
\sum_{n=0}^{\infty}F_0(4n+1)q^n
=
\dfrac{f_2^{13}}{f_1^{12}f_4^2}.
\end{equation}
Employing \eqref{2t} in \eqref{e26},
we obtain
\begin{equation}\label{phps}
\sum_{n=0}^{\infty}F_0(4n+1)q^n
\equiv
\dfrac{f_2f_4^2}{f_1^4}
\equiv\dfrac{f_4^2}{f_2}
\left(\dfrac{f_2}{f_1^2}\right)^2
\equiv\psi(q^2)\phi(q)^2
\pmod 8.
\end{equation}
Using \eqref{psi} and \eqref{gp} in \eqref{phps}, we obtain
\begin{equation}\label{fg8}
\sum_{n=0}^{\infty}F_0(4n+1)q^n
\equiv
\left(\sum_{a=0}^{\infty}q^{a(a+1)}\right)
\left(\sum_{m=0}^{\infty}r_2(m)q^m\right) 
\equiv
\sum_{n=0}^{\infty}
\left(
\sum_{\substack{a\geq0\\a(a+1)\leq n}}
r_2\bigl(n-a(a+1)\bigr)
\right)q^n.
\end{equation}
By  Fundamental Theorem of Arithmetic (FTA), any integer $m\geq1$ can be expressed as $m=2^\alpha M,$ where $2\nmid M$ and $\alpha\geq0.$ Thus, from \eqref{r_2}, we obtain
\begin{equation}\label{fg8e}
r_2(m)
=
4\sum_{\substack{d\mid m\\2\nmid d}}(-1)^{(d-1)/2}
=
4\sum_{d\mid M}(-1)^{(d-1)/2} 
\equiv
4\sum_{d\mid M}1
\equiv
4d(M)
\pmod 8,
\end{equation}
where $d(M)$ is the number of positive divisors of $M.$ Again, if $p_i$ are distinct primes  and $\alpha_i\geq1$ are integers for ($i=1,\;2,\;\ldots$), then by FTA, 
$M=p_1^{\alpha_1}p_2^{\alpha_2}\cdots p_k^{\alpha_k}$
and 
\begin{equation}\label{fg10}
d(M)
=
\prod_{j=1}^{k}(\alpha_j+1).
\end{equation}
Employing \eqref{fg10} in \eqref{fg8e}, we obtain 
\begin{equation}\label{fg11}
r_2(m)\equiv
\begin{cases}
4 \pmod 8, & \text{if the odd part of }m\text{ is a perfect square},\\
0 \pmod 8, & \text{otherwise}.
\end{cases}
\end{equation}
Employing \eqref{fg11} in \eqref{fg8} and further comparing coefficients of like powers of $q$, we obtain
\begin{equation}\label{fg12}
F_0(4n+1)
\equiv
\varepsilon(n)+4~|\mathcal{E}(n)|
\pmod 8,
\end{equation}
where
$
\mathcal{E}(n)=
\left\{
a\geq0:
n-a(a+1)=2^\alpha x^2
\text{ for some integer } \alpha\geq0,\ x\geq1
\right\}
$ and $|\cdot|$ denotes the order of the set. 

Define, 
\begin{align}
\label{fg15}\mathcal{E}_1(n)
&=
\left\{
a\geq0:
n-a(a+1)=y^2
\text{ for some integer } y\geq1
\right\},\\
\intertext{and}
\label{fg16}\mathcal{E}_2(n)
&=
\left\{
a\geq0:
n-a(a+1)=2y^2
\text{ for some integer } y\geq1
\right\}.
\end{align}
Clearly, 
\begin{equation}\label{fg17}
|\mathcal{E}(n)|=|\mathcal{E}_1(n)|+|\mathcal{E}_2(n)|.
\end{equation}
Again,
$$r_2(4n+1)
=|\{(u,v)\in\mathbb Z^2:u^2+v^2=4n+1\}|.$$
This implies, for $(u,v)\in \{(u,v)\in\mathbb Z^2:u^2+v^2=4n+1\},$ 
both $u$ and $v$ cannot be even at the same time; that is, one of   $u$ and $v$ is even and the other is  odd. If $u=2a+1$ and $v=2y,$ for nonnegative integers $a,y\geq0$,  then 
\begin{equation}\label{ha3}
n-a(a+1)=y^2\iff
(2a+1)^2+(2y)^2=4n+1 .
\end{equation}
If $y=0,$ then by  \eqref{ha3},  $(u,v)=(\pm(2a+1),0),(0,\pm(2a+1))$; and if  $y\geq1,$ then  $(u,v)=(\pm(2a+1),\pm 2y),(\pm 2y,\pm(2a+1)).$
Therefore,
\begin{equation}\label{ha8}
r_2(4n+1)
=
4\varepsilon(n)+8|\mathcal E_1(n)|.
\end{equation}
Employing \eqref{r_2.4n+1} in \eqref{ha8}, we obtain
\begin{equation}
    |\mathcal E_1(n)|
=
\dfrac{1}{2}
\left(
\sum_{d\mid 4n+1}\left(\dfrac{-1}{d}\right)
-
\varepsilon(n)
\right).
\label{fg21}
\end{equation}
If we consider  $\mathcal{E}_2(n),$ then 
\begin{equation}\label{fg22}
n-a(a+1)=2y^2
\Longleftrightarrow
4n+1=(2a+1)^2+8y^2.
\end{equation}
Let $r_{1,2}(N)$ denote the number of representations of $N$ in the form $u^2+2v^2,$ where $(u,v)\in\mathbb{Z}^2.$\\ Then
\begin{equation}\label{fg23}
r_{1,2}(4n+1)=2\varepsilon(n)+4|\mathcal{E}_2(n)|.
\end{equation}
Employing \eqref{B3} in \eqref{fg23}, we obtain
\begin{equation}\label{fg25}
|\mathcal{E}_2(n)|
=
\dfrac{1}{2}
\left(
\sum_{d\mid 4n+1}\left(\dfrac{-2}{d}\right)
-
\varepsilon(n)
\right).
\end{equation}
Employing \eqref{fg21} and \eqref{fg25} in  \eqref{fg17},  we obtain
\begin{align}
|\mathcal{E}(n)|
&=
\dfrac{1}{2}
\left(
\sum_{d\mid 4n+1}\left(\dfrac{-1}{d}\right)
-
\varepsilon(n)
\right)
+
\dfrac{1}{2}
\left(
\sum_{d\mid 4n+1}\left(\dfrac{-2}{d}\right)
-
\varepsilon(n)
\right) \notag\\
&=
\dfrac{1}{2}
\left(
\sum_{d\mid 4n+1}\left(\dfrac{-1}{d}\right)
+
\sum_{d\mid 4n+1}\left(\dfrac{-2}{d}\right)
-
2\varepsilon(n)
\right) \notag\\
&=
\dfrac{1}{2}\bigl(\mathcal{D}(n)-2\varepsilon(n)\bigr).
\label{fg26}
\end{align}
Employing \eqref{fg26} in \eqref{fg12}, we obtain
\begin{equation}
F_0(4n+1)
\equiv
\varepsilon(n)+2\bigl(\mathcal{D}(n)-2\varepsilon(n)\bigr)
\pmod 8.
\end{equation}
This completes the proof of \eqref{fg3}. \\\\
Again,  from \eqref{fg26}, we note that
\begin{equation}
\mathcal{D}(n)-2\varepsilon(n)=2|\mathcal{E}(n)|,\notag
\end{equation}
Then the following four cases are possible:
\begin{align}
\varepsilon(n)=0,\quad \mathcal{D}(n)\equiv0\pmod4
&\Longrightarrow
F_0(4n+1)\equiv0\pmod8,\notag\\
\varepsilon(n)=0,\quad \mathcal{D}(n)\equiv2\pmod4
&\Longrightarrow
F_0(4n+1)\equiv4\pmod8,\notag\\
\varepsilon(n)=1,\quad \mathcal{D}(n)\equiv2\pmod4
&\Longrightarrow
F_0(4n+1)\equiv1\pmod8,\notag\\
\varepsilon(n)=1,\quad \mathcal{D}(n)\equiv0\pmod4
&\Longrightarrow
F_0(4n+1)\equiv5\pmod8.\notag
\end{align}
This proves \eqref{fg4} for the case $i=0$. It is easy to see that $F_0(2n+1)=F_1(2n+1),$ so the case  $i=1$ of \eqref{fg4} follows trivially.
\end{proof}

\end{document}
